\chapter{Los datos: de dónde sale cada cifra}
\label{cap:datos}

\section{La regla: solo fuentes oficiales, descargadas con código}
Todas las fuentes se descargan con programas (nunca a mano), a \codigo{datos/bronze/<fuente>/<fecha>/}, \textbf{sin
modificarlas}. Cada archivo queda registrado en \codigo{datos/bronze/MANIFIESTO.csv} con su huella SHA-256, su tamaño, su
dirección de origen, la fecha de descarga y su número de registros. Resultado de la Fase 1: \textbf{353 archivos,
8,134,802 registros y 824 MB}.

Dos razones para hacerlo así:
\begin{itemize}
\item \textbf{Reproducible}: cualquiera puede volver a correr la descarga y obtener lo mismo, o ver qué cambió.
\item \textbf{Rastreable}: cualquier cifra del proyecto se puede seguir hasta el archivo oficial del que salió
  (\codigo{docs/trazabilidad.md}).
\end{itemize}
Lo que \textbf{no} se usa: TripAdvisor y Google Maps, porque sus términos prohíben extraer sus datos; el semáforo de
sargazo del estado, porque solo se publica en Facebook, cuyos términos también lo prohíben.

\section{Las fuentes, una por una}

\subsection*{D1 · SITUR-Q (sistema estatal de información turística)}
\begin{itemize}
\item \textbf{Qué es}: el sistema del gobierno de Quintana Roo; publica indicadores por destino y mes.
\item \textbf{Cómo se descargó}: su página consulta una API con una llave pública incrustada. Al principio rechazaba las
  peticiones (``Origin not allowed''); se resolvió enviando los mismos encabezados que envía el navegador.
\item \textbf{Qué trae}: ocupación hotelera, cuartos y hoteles, Tren Maya (subidas y bajadas por estación), cruceristas,
  cruces desde Belice, llegadas por avión, visitas a zonas y afluencia. 6,607 registros.
\item \textbf{Hallazgos}: \emph{no publica ocupación hotelera de 2025–2026 para ningún destino} (viene en ceros); las
  llegadas por avión de 2025 vienen en cero en los cuatro aeropuertos; el indicador ``Turista – Afluencia'' devolvió
  error 500 en las 120 consultas. Todo eso se declaró como hueco.
\item \textbf{Para qué}: Radar del sur, pronóstico de Chetumal (cruces de Belice), conciliación con DataTur.
\end{itemize}

\subsection*{D2 · DataTur, ocupación hotelera (SECTUR)}
\begin{itemize}
\item \textbf{Qué es}: el monitoreo hotelero federal. 135 archivos semanales y 31 mensuales.
\item \textbf{Qué trae}: cuartos disponibles, cuartos ocupados y ocupación de 66 centros del país; los 7 de Quintana Roo
  (Cancún, Riviera Maya, Playa del Carmen, Playacar, Akumal, Cozumel e Isla Mujeres) tienen \textbf{239 semanas}, de enero
  de 2022 a julio de 2026 (cada archivo trae también los dos años anteriores).
\item \textbf{Hallazgo}: hubo 2,804 cifras corregidas entre una publicación y otra; se conserva la más reciente.
\item \textbf{Para qué}: la única ocupación medida de 2025–2026; Radar del norte, cadena de Markov y Torre en vivo.
\end{itemize}

\subsection*{D3 · DataTur, llegadas por nacionalidad (Unidad de Política Migratoria)}
521,363 renglones de 2012 a julio de 2026: extranjeros que llegan en avión, por aeropuerto, país y sexo. Solo cuenta
extranjeros (no hay país ``México''). Para la campaña: de dónde viene el visitante que aterriza en Cancún y cuántos
llegan a Chetumal (250 en 2025).

\subsection*{D4 · DataTur, INAH, AFAC, cruceros y Compendio}
\begin{itemize}
\item \textbf{BdINAH}: visitantes por zona arqueológica, mes y tipo (nacional o extranjero), 2016–2026. Es la serie
  \textbf{medida} que sí llega a 2026 en el sur, y la base del pronóstico.
\item \textbf{AFAC}: pasajeros por aerolínea en el país. \textbf{Cruceros}: arribos y pasajeros por puerto.
\item \textbf{Compendio 2024}: cuartos-noche por categoría de hotel (estrellas) en 70 centros.
\item En total, 390,228 registros.
\end{itemize}

\subsection*{D5 · Rest-Mex 2025 (licencia CC-BY-4.0)}
208,051 reseñas de turistas en México con calificación de 1 a 5. En Quintana Roo, 85,993 (85,987 sin duplicados), solo
de Tulum, Isla Mujeres y Bacalar. Para la minería de texto de la campaña.

\subsection*{D6 · DENUE del INEGI (los 32 estados)}
6,138,075 negocios con giro, tamaño y coordenadas. Se procesó el país completo, y no solo Quintana Roo, por dos razones:
la rúbrica pide volumen real, y el total nacional da contexto (14.3\,\% de los negocios del país son turísticos; en
Quintana Roo, 19.7\,\%). Para la oferta turística, los criterios y el planeador.

\subsection*{D7 · Censo 2020, ITER del INEGI}
Población y viviendas de 2,243 localidades de Quintana Roo. Se usa \textbf{por localidad}, no por municipio, porque
cuatro de los cinco lugares están en el mismo municipio (Othón P. Blanco) y con la cifra municipal tendrían el mismo
denominador.

\subsection*{D8 · Open-Meteo (reanálisis ERA5)}
Clima de 8 puntos: diario desde 1950 (224,232 días) y horario desde 2019 (542,784 horas), sin huecos. Un
\emph{reanálisis} es un modelo meteorológico que combina observaciones reales; no es una estación, y así se declara. Para
la regresión con clima y para detectar clima raro.

\subsection*{D9 · HURDAT2 de la NOAA}
Todas las tormentas del Atlántico de 1851 a 2025: 55,524 posiciones de 1,988 tormentas, cada seis horas. Para la
probabilidad de tormenta.

\subsection*{D10 · FRED (Reserva Federal de St. Louis)}
Pesos por dólar de cada día hábil desde 1993 (8,575 días) e inflación de EE.\,UU. Para la sensibilidad al tipo de cambio
y para pasar a pesos los costos de los anuncios.

\subsection*{D11 a D15}
ENDUTIH 2025 (uso de internet y aplicaciones, PDF oficial del INEGI), los polígonos de los municipios, los costos por
clic y conversiones de anuncios de viajes (WordStream y LocaliQ, promedios de EE.\,UU. usados como supuesto etiquetado) y
las fotos con licencia libre de Wikimedia Commons.

\section{Lo que no existe (huecos declarados)}
\begin{center}
\begin{tabularx}{\linewidth}{lY}
\encabezado \h{Hueco} & \h{Qué no se puede afirmar por su culpa} \\
Ocupación hotelera del sur 2025–2026 & Cuántos cuartos se llenan hoy en Chetumal; por eso el sur se mide con llegadas reales \\
Estado de origen, edad e ingreso del visitante nacional & Una persona nacional más detallada \\
Conversión de Facebook para turismo & El rendimiento exacto de Facebook; se prueba con 3 valores \\
Tormentas de 2026 & Si hubo tormentas este año; la Torre dice ``sin dato'' \\
Reseñas de los cinco lugares & Qué opinan los turistas del sur; las reseñas son de destinos llenos \\
Derrama económica por lugar & Cuánto dinero deja la campaña; la de SITUR-Q no trae unidad y termina en 2024 \\
\bottomrule
\end{tabularx}
\end{center}
